\documentclass[10pt,a4paper]{article}
\usepackage[margin=1.55cm]{geometry}
\usepackage{xcolor}
\usepackage{amsmath,amssymb}
\usepackage{graphicx}
\usepackage{enumitem}
\usepackage{fontspec}
\usepackage[CJKmath=true]{xeCJK}
% CJK main font is resolved at COMPILE time on the actual machine, not baked in at
% generation time, so a .tex generated on one OS still renders on another (the older
% Python-side probe hardcoded one OS's font name into the file and broke when moved).
% Walk a cross-platform priority list and use the first font that exists; AutoFakeBold
% + AutoFakeSlant give bold/italic CJK even on faces that ship no bold/italic cut.
\newcommand{\setCJKto}[1]{\setCJKmainfont{#1}[AutoFakeBold=true,AutoFakeSlant=0.2]}
\IfFontExistsTF{Noto Sans CJK SC}{\setCJKto{Noto Sans CJK SC}}{%        % Linux (fonts-noto-cjk); cross-platform if installed
 \IfFontExistsTF{Source Han Sans SC}{\setCJKto{Source Han Sans SC}}{%   % Adobe 思源黑体; cross-platform
  \IfFontExistsTF{PingFang SC}{\setCJKto{PingFang SC}}{%               % macOS 10.11+
   \IfFontExistsTF{Hiragino Sans GB}{\setCJKto{Hiragino Sans GB}}{%    % older macOS
    \IfFontExistsTF{Microsoft YaHei}{\setCJKto{Microsoft YaHei}}{%     % Windows
     \IfFontExistsTF{WenQuanYi Micro Hei}{\setCJKto{WenQuanYi Micro Hei}}{% % older Linux
      \setCJKto{Songti SC}}}}}}}                                       % macOS bottom fallback
\definecolor{cblue}{HTML}{4C72B0}
\setlength{\parskip}{3pt}\setlength{\parindent}{0pt}
\setlist[enumerate]{topsep=2pt,itemsep=2pt,parsep=1pt,partopsep=0pt}
\setlist[itemize]{topsep=2pt,itemsep=2pt,parsep=1pt,partopsep=0pt}
% Let prose-embedded inline math break across lines and absorb small overflows, so simple
% formulas inserted mid-sentence stop running past the right margin.
\tolerance=1200
\emergencystretch=2.5em
\relpenalty=20
\binoppenalty=40
\hfuzz=1pt
\begin{document}

\section*{量化 World Action Models 的因果后果校准}
\textbf{方法名称：} One-Transition Consequence PTQ (OTC-PTQ)
\section*{研究动机}
对于一个耦合 future-video/action 的 World Action Model（WAM），post-training quantization（PTQ）校准不只是复现局部 tensor 数值。一次 numerical intervention 可能改变模型预测的完整 action chunk；只执行第一个 action 后，simulator 的 next observation 也可能改变。实际瓶颈是选择哪些 weight/activation（W/A）site 保持 high precision，因为 local PTQ discrepancy 并不能稳定识别 contact-transition states 上的 consequence。因此，OTC-PTQ 把完整 action chunk 当作中间对象，把随后产生的一次 control transition 作为可测量边界。

当前 Phase 0 frontier 使这个边界可以被测试。QuantWAMs（arxiv:2607.28405v1）明确处理 joint video/action sensitivity 和 closed-loop protection，QuantWM（arxiv:2602.02110v1）提供不同的 latent rollout/planning sensitivity 分支，Fast-WAM 则提供名为 idm 的 feasibility path，但没有提供 numerical PTQ consequence signal。结合已经说明的 RoboTwin 2.0/LIBERO manipulation split 和 Fast-WAM-only staged validation budget，现在可以进行一个小规模 one-transition diagnostic；前提是 pre-run manifest 中固定 exact Optional IDM checkpoint identity。

几个相邻方向已经有工作覆盖。QuantWAMs 测量 joint video/action sensitivity，并修复 closed-loop protection schedule；QuantWM 研究 weight/activation PTQ sensitivity 和 long-horizon latent planning；Fast-WAM 没有提供用于 one-site W/A protection allocation 的 numerical PTQ consequence signal。更早的 task-consequence 工作------VAML、Model Advantage / Value-Aware Model Learning 和 Deep Task-Based Quantization------已经说明 downstream consequence 可以指导 targeting；SQIL 则通过 saliency-weighted QAT 和 action distillation 覆盖 critical-state protection。因此，更窄的 gap 是：当其他 target site 全部保持 FP16 时，一次 numerical W/A intervention 能否根据 paired full-action-chunk 和 next-observation discrepancy，提供可复现的 one-control-transition proxy。OTC-PTQ 不把这些更广泛的 downstream-consequence 或 critical-state-targeting 原则宣称为新贡献。

如果这个 gap 被填补，领域将得到一个有界 diagnostic，用来识别 local PTQ discrepancy 何时不再预测 contact states 上的一次 control transition consequence。它补充 QuantWAMs 的 closed-loop protection 分支和 QuantWM 的 latent-planning sensitivity 分支。它还可以检验 \(C_{r}-ranked\) schedule 是否能在相同 precision budget 下迁移到 held-out contact-sensitive failures 和 next-observation deviation；permutation 可以把 consequence association 与 quantizer 和 site-count controls 区分开来。这不会证明 task-loss identity、full-horizon/replan marginal、joint-low-bit optimality 或 bytes/latency objective。
\section*{方法}
\subsection*{M1\_background}
\emph{提供可 replay 的 calibration state，并固定 paired measurement 所需的 exact version 与 configuration boundary。}
\begin{enumerate}[leftmargin=*,start=1]
  \item 先冻结 manifest。 在任何 paired measurement 之前，为每个 state 建立不可变 calibration record，记录 state、checkpoint/configuration identities、controller/RNG/history/cache restoration contract 与 environment identity。 manifest 中必须提供 Fast-WAM 与 Optional IDM 的 exact checkpoint identities 和 exact-match fields；缺失或不匹配时 halt。【作者需决定：补齐不可变 checkpoint identities、revision/digest 与 exact-match fields；不得 fallback 到其他 artifact。】 environment record 还必须固定 CPU thread/process limits、accelerator model、visible-memory class 与 memory-allocation policy。【作者需决定：在 manifest 中补齐这些 execution-environment fields，并明确 A100-40GB 与 80GB-class 边界。】 任何必需 identity 或 restoration check 缺失或不匹配，都必须在建立 paired branches 之前 halt。
\par\emph{为什么：} 只有两个 branch 都从同一个可恢复 state 开始，paired consequence 才能解释。这样可以保留 version boundary，不会把未经验证的 checkpoint 当作可互换对象。
\end{enumerate}
\subsection*{M2\_one\_transition\_proxy}
\emph{建立有界的 single-site W/A intervention，并计算 one-control-transition \(C_{r}\) signal。}
\begin{enumerate}[leftmargin=*,start=2]
  \item 建立一个 versioned site registry，给 weight/activation (W/A) site 配置稳定的 \(site_{id}\)、module path 或 graph edge、tensor shape、dtype、capture point 与 quantization granularity。 对每个 candidate，先恢复 第1步（先冻结 manifest） baseline，再应用只包含一个 \(site_{id}\) 的 mask；在注册的 graph edge 上量化选定的 weight tensor 或截取选定的 activation，其余 target sites 全部保持 16-bit floating-point (FP16)。 记录 operator parameters，并验证 mask 只改变了被选中的 site。 必须把 low-bit operator 声明为 fake quantization（图级仿真）或真实的 low-bit execution kernel，并在 calibration 与 held-out evaluation 中保持同一条已声明路径。 【作者需决定：选择并登记 \(b_{lo}\)、signedness、rounding、clipping、scale/zero-point、granularity、activation calibration data，以及 fake-quantization 与 real-kernel mode；若声称使用真实 kernel，还要记录 backend/version 并做 runtime verification。】 这只是 single-site、FP16-background 的 intervention，不是 joint-low-bit marginal 或 joint-optimal schedule；QuantWM 与 QuantWAMs 的更广 quantization scope 仍属于 prior-art boundary。
\noindent\emph{只在 numerical site r 应用 low-bit operator，其余目标 site 保持 FP16。}
\par\noindent{\small\ttfamily [eq~1, raw LaTeX, not typeset] \textbackslash\{\}mathcal\{Z\}\textasciicircum{}\{(r)\}\_j=\textbackslash\{\}begin\{cases\}Q\_\{b\_\{\textbackslash\{\}mathrm\{lo\}\}\}(\textbackslash\{\}mathcal\{Z\}\_r),\&j=r\textbackslash\{\}\textbackslash\{\}\textbackslash\{\}mathcal\{Z\}\textasciicircum{}\{\textbackslash\{\}mathrm\{FP16\}\}\_j,\&j\textbackslash\{\}ne r\textbackslash\{\}end\{cases\}}
\par\emph{为什么：} 这会隔离有界的 single-site intervention，也避免把 joint coupling、joint-low-bit marginal effects 或 joint optimality 当作新的 claim。
  \item 从匹配的 restored states 分别运行 reference branch 与 single-site branch。执行任何 action 之前保存完整 action chunks，包括 time length、action ordering、valid-step mask 与 dtype；每个 branch 分别恢复同一 state，只执行第一个 action，并保存 immediate next observation 以及 terminal 或 simulator-error flags。定义 \(d_{a}(j,r)\)：完整 action chunk 的 MSE 除以 \(max(sigma_{a}^{2}, epsilon_{a}^{2})\)，其中 \(sigma_{a}^{2}\) 是 \(D_{cal}\) 中 FP16 full action chunks 所有 scalar coordinates 的 population variance；定义 \(d_{o}(j,r)\)：next observation 的 mean per-pixel absolute difference 除以 \(max(sigma_{o}, epsilon_{o})\)，其中 \(sigma_{o}\) 是 calibration 中 per-pixel standard deviations 的 mean。使用 normalized units 中的 \(epsilon_{a}=epsilon_{o}=1e-3\)，默认 \(lambda_{a}=lambda_{o}=1\)，预注册 \(lambda_{a}/lambda_{o}\) sweep \({0.5,1,2}\)，并将 \(C_{r}\) 定义为 \(lambda_{a} d_{a}+lambda_{o} d_{o}\) 的 episode mean，在完整 episode 上计算 bootstrap intervals。定义 open-loop local surrogate \(L_{jr}\) 为 paired hook outputs 所有 elements 的 MSE，\(L_{r}\) 为其 episode mean。仍需明确 valid-element mask、observation preprocessing，以及对 terminal、invalid、simulator-error、missing-observation 与 variable-length rows 一致适用的 inclusion policy。\(C_{r}\) 仍是 one-control-transition proxy，不是 task loss、后续 rollout/replan marginal 或 jointly-low-bit marginal。 pixel term 的 valid-element mask 与 observation preprocessing 仍需补齐。【作者需决定：在 held-out evaluation 前记录两者，并对两个 branch 一致执行。】terminal、invalid、simulator-error、missing-observation 与 variable-length rows 也需要同一套 inclusion policy。【作者需决定：说明每类 row 是 retained、masked 还是 excluded，并记录每个 exclusion reason。】
\noindent\emph{\(d_{a}\) 是完整 action chunk 的 MSE 除以 \(max(sigma_{a}^{2}, epsilon_{a}^{2})\)，\(d_{o}\) 是 next observation 的 per-pixel 绝对差除以 \(max(sigma_{o}, epsilon_{o})\)。在 normalized units 中使用 \(epsilon_{a}=epsilon_{o}=1e-3\)，默认 \(lambda_{a}=lambda_{o}=1\) 并执行预注册 sensitivity sweeps；\(C_{r}\) 按 episode mean 聚合，并在完整 episode 上计算 bootstrap intervals。}
\par\noindent{\small\ttfamily [eq~2, raw LaTeX, not typeset] d\_a(j,r)=\textbackslash\{\}frac\{\textbackslash\{\}operatorname\{MSE\}(\textbackslash\{\}mathbf a\_\{jr\},\textbackslash\{\}mathbf a\_\{j0\})\}\{\textbackslash\{\}max(\textbackslash\{\}sigma\_a\textasciicircum{}2,\textbackslash\{\}epsilon\_a\textasciicircum{}2)\},\textbackslash\{\}quad d\_o(j,r)=\textbackslash\{\}frac\{\textbackslash\{\}operatorname\{mean\}\_\{\textbackslash\{\}mathrm\{pixel\}\}\textbackslash\{\}left|\textbackslash\{\}mathbf o\_\{(j+1)r\}-\textbackslash\{\}mathbf o\_\{(j+1)0\}\textbackslash\{\}right|\}\{\textbackslash\{\}max(\textbackslash\{\}sigma\_o,\textbackslash\{\}epsilon\_o)\},\textbackslash\{\}quad C\_r=\textbackslash\{\}operatorname\{mean\}\_\{j\textbackslash\{\}in D\_\{\textbackslash\{\}mathrm\{cal\}\}\}\textbackslash\{\}left(\textbackslash\{\}lambda\_a d\_a(j,r)+\textbackslash\{\}lambda\_o d\_o(j,r)\textbackslash\{\}right)}
\par\emph{为什么：} 这定义了一个表示 action-mediated consequence 的 one-control-transition proxy。\(C_{r}\) 不是 task loss、full-horizon/replan marginal，也不是 jointly-low-bit marginal。
\end{enumerate}
\subsection*{M3\_matched\_protection}
\emph{使用 \(C_{r}\) 和 local discrepancy 形成 matched protection schedule，同时分开表达不同 budget 的含义。}
\begin{enumerate}[leftmargin=*,start=4]
  \item 使用同一个 site registry 建立两份 deterministic protection schedule：一份按 \(C_{r}\) 排序，另一份按 local discrepancy 排序。将 \(B_{site}\) 定义为 matched local-fidelity baseline 在 preregistered run manifest 中记录的 nonnegative integer high-precision-site count，并在本方法中固定 B 等价于 \(B_{site}\)。按 score 降序排列，ties 按 registry order 处理；把恰好 \(top-B_{site}\) sites 保持为 FP16，其余 target sites 保持为 \(b_{lo}\)。\(B_{site}\) 的数值和 manifest entry 仍是 open inputs；如果 entry 缺失或 \(B_{site}\) 超过 registry 中的 site 数量，就 halt。persistent weight bytes（包括 copies 与 scales）、real kernel 下的 peak live activation workspace、scale/copy overhead 与 measured latency 必须另外记录，不能替代 site-count budget。 numeric \(B_{site}\) manifest entry 仍是 open。【作者需决定：提供 nonnegative integer baseline site count；缺失或超过 registry size 时 halt。】deployment bytes、activation workspace、scale/copy overhead 与 latency 仍需独立 measurement protocol。【作者需决定：提供 packing/timing protocol，或明确把这些 outcome 标为 out of scope。】
\noindent\emph{按 \(C_{r}\) 选择前 \(B_{site}\) 个 numerical site；site count 不等同于 bytes 或 latency 目标。}
\par\noindent{\small\ttfamily [eq~3, raw LaTeX, not typeset] \textbackslash\{\}mathcal\{P\}\_\{C\_r\}(B\_\{\textbackslash\{\}mathrm\{site\}\})=\textbackslash\{\}operatorname\{TopB\}\_\{\textbackslash\{\}mathrm\{site\}\}(\textbackslash\{\}\{(r,C\_r):r\textbackslash\{\}in\textbackslash\{\}mathcal\{R\}\textbackslash\{\}\})}
\par\emph{为什么：} 这个比较检验 proxy 是否会在匹配的 precision budget 下改变 protection choice，而不是引入 storage 或 runtime objective。
\end{enumerate}
\subsection*{M4\_validation}
\emph{通过 permutation 和 oracle controls 检验 contact-transition divergence 与有界 transfer。}
\begin{enumerate}[leftmargin=*,start=5]
  \item 使用不可变的 第1步（先冻结 manifest） manifest，在命名的 held-out RoboTwin 2.0/LIBERO manipulation split 上运行带有 sigma\_shift\_FW=1.0 的 idm entrypoint 与 infer\_action。 对两份 schedule 使用相同的声明性 \(b_{lo}\)、B、\(B_{site}\)、preprocessing、reset procedure、controller restoration 与 branch order。 在查看结果之前，列出 held-out task identifiers、episode 与 initial-state identifiers、seeds、split membership、来自 simulator events 的 contact predicate，以及 task success/failure rule。 【作者需决定：提供 exact split 与 contact/failure definitions，包括 contact 是否必须在第一个 action 之前发生，还是允许在之后发生。】 对每个 paired episode，为两份 schedule 恢复同一个 initial state，运行声明的 evaluation horizon，并记录 task outcome 与带有 aggregation window 的 next-observation deviation。 【作者需决定：固定 idm 与 infer\_action entrypoint revisions 和 arguments、action horizon、termination rule，并决定 next-observation deviation 只看 first transition，还是覆盖 later transition/full episode；后者必须与 \(C_{r}\) 分开汇总。】 建立 permutation control：在相同 site identifiers 之间重新分配已记录的 \(C_{r}\) values，同时保持 score multiset、B、\(B_{site}\)、tie-breaking 与其他 schedule rules 不变。 指定的 oracle simulator-state branch 只能作为标注清楚的 positive-control branch 运行。 【作者需决定：指定 permutation seed 与 replicate/list、oracle state source、branch insertion point 以及它所控制的 comparison；不能用 oracle 替代真实 association。】 保持解释边界：\(C_{r}\) 仍是 one-control-transition ranking proxy；task failures 与 full-horizon summaries 是独立的 downstream outcomes，不是 task loss，也不是 full-horizon/replan marginal。 QuantWM 与 QuantWAMs 继续保留其更广的 world action model (WAM) quantization boundaries；VAML、Model Advantage / Value-Aware Model Learning、Deep Task-Based Quantization 与 SQIL 继续代表 broad consequence、value-aware、task-based 或 critical-state-protection principles 的 prior-art boundary。
\noindent\emph{按 state subset 比较 local score 与 \(C_{r}\) 的 rank correlation，预期差异集中在 contact-transition states。}
\par\noindent{\small\ttfamily [eq~4, raw LaTeX, not typeset] \textbackslash\{\}Delta\textbackslash\{\}rho=\textbackslash\{\}rho\_\{\textbackslash\{\}mathrm\{noncontact\}\}-\textbackslash\{\}rho\_\{\textbackslash\{\}mathrm\{contact\}\},\textbackslash\{\}quad \textbackslash\{\}rho\_s=\textbackslash\{\}operatorname\{Spearman\}(\textbackslash\{\}\{L\_r\textbackslash\{\}\}\_s,\textbackslash\{\}\{C\_r\textbackslash\{\}\}\_s)}
\par\emph{为什么：} contact-transition subset 和 permutation 是 mechanism tests：candidate 预测真实 \(C_{r}\) association 会带来 ranking divergence 和 transfer，而 permutation 后这个 advantage 会消失。oracle branch 用来检查 calibration noise。
\end{enumerate}
\end{document}
